\begin{corollary}[Two-phase mixing principle on the Pareto front: asymptotically optimal fiber structures with at most two atomic lengths]\label{cor:pom-fiber-ab-two-phase-mixing}
Adopt the notation of Theorem \ref{thm:pom-fiber-ab-convex-hull}, and let \(K:=\overline{\mathrm{conv}}(\mathcal P)=\mathcal R\subset\RR^2\).
Then for any support direction \(u=(\alpha,\beta)\neq(0,0)\), its exposed face
\[
F(u):=\left\{x\in K:\ \langle u,x\rangle=\sup_{y\in K}\langle u,y\rangle\right\}
\]
can only be a point or a line segment.
Therefore any exposed point \(x\in F(u)\) on a Pareto front belongs to a closed segment spanned by two points \(x_-,x_+\in F(u)\):
\[
x=\theta x_-+(1-\theta)x_+,\qquad \theta\in[0,1].
\]
In particular, when the endpoints can be taken as atomic points \(x_\pm=p_{\ell_\pm}\in\mathcal P\), the optimal decomposition in the scale limit requires at most a mixture of two atomic lengths \(\ell_-,\ell_+\in\NN\) (single-phase is the degenerate case \(\ell_-=\ell_+\)).
When the support direction varies continuously, switching of endpoint pairs can only occur at the enumerable threshold set generated by pairwise slopes of atomic points, thereby inducing a family of discrete ``phase boundaries.''
\end{corollary}
\begin{proof}
Since \(F(u)\subset\{x\in\RR^2:\langle u,x\rangle=\mathrm{const}\}\) and \(u\neq 0\), this support hyperplane is a line in \(\RR^2\).
Also, \(F(u)\) is a convex compact set, so it can only be a point or a segment; in the segment case, any point can be written as a convex combination of the two endpoints.
Finally, if the segment endpoints are given by two atomic points \(p_{\ell_-},p_{\ell_+}\), the slope of the support direction is determined by the slope of the line joining this pair; hence the set of directions where switching can occur is contained in the set of all pairwise slopes, thus is enumerable.
\qed
\end{proof}

